\documentclass{beamer}

\usepackage[utf8]{inputenc}
\input{../helper}

\title{Introdução à Computação \\ passagem de valores por referência}
\author{Alexandre Rademaker}
\date{}

\begin{document}

\begin{frame}
 \maketitle
\end{frame}


\begin{frame}[fragile,allowframebreaks]{troca (swap)}

  Queremos trocar dois valores de lugar. 

  \lstinputlisting[style=normalc,captionpos=t]{swap0.c}
  
  \framebreak

  Como vimos, C usa como padrão a estratégia de chamar as funções por
  valor
  (\href{https://en.wikipedia.org/wiki/Evaluation_strategy#Call_by_value}{call-by-value}).

  A função \ilcode{swap} recebeu uma cópia dos valores, isto é, os
  valores foram copiados para uma nova região da memória.

  Mas ponteiros nos permitem controlar esta estratégia.

  \framebreak

  \lstinputlisting[style=normalc,captionpos=t,caption=swap1.c]{swap1.c}

\end{frame}

\begin{frame}{organização da memória}

  O que aconteceu na memória com swap?

  \foreach \n in {2,...,37}{
    \includegraphics<\n>[width=.9\textwidth]{swap/swap-\n.png}}

\end{frame}


\begin{frame}{Exemplos extras}
  Arquivos em \texttt{src/} relacionados ao tema desta semana, para
  estudo complementar (não foram apresentados em aula):
  \begin{itemize}
    \item \texttt{arrays.c}
    \item \texttt{ascii0.c}, \texttt{ascii1.c}
    \item \texttt{buggy2.c}
    \item \texttt{scores2.c}, \texttt{scores3.c}, \texttt{scores4.c}
  \end{itemize}
\end{frame}

\end{document}

%%% Local Variables:
%%% mode: latex
%%% TeX-master: t
%%% End:
